% ============================================================================
% Fichier  : partie4/P04_C12_itil_systeme_valeur.tex
% Titre    : ITIL 4 : le Système de Valeur des Services
% Auteur   : MINKA MI NGUIDJOÏ Thierry Emmanuel
% Version  : 1.0 -- octobre 2026
% ============================================================================
% Décision de l'auteur (29/10/2026) : la Partie IV s'appuie sur ITIL 4,
% avec une note sur ITIL version 5 (PeopleCert, 2026). Sources publiques
% uniquement : les définitions sont reformulées, aucune n'est citée
% textuellement.
% ============================================================================

\chapter{ITIL 4 : le Système de Valeur des Services}
\label{chap:itil-svs}
\chaptermark{ITIL 4 : le système de valeur}

\epigraph{%
    « La valeur ne se livre pas.\\
    Elle se crée avec celui qui reçoit le service. »%
}{Idée centrale d'ITIL~4}

\niveauChapitre{Intermédiaire}{%
    Maîtrise du Chapitre~\ref{chap:cobit-archi} (architecture de COBIT~2019)
    et du Chapitre~\ref{chap:referentiels} (place d'ITIL parmi les
    référentiels).%
}

\begin{boxobjectifs}
\begin{itemize}
    \item Retracer l'évolution d'ITIL (v2, v3 et son cycle de vie, ITIL~4).
    \item Définir service, valeur, produit et relation de service.
    \item Décrire le système de valeur des services et sa chaîne de valeur.
    \item Appliquer les sept principes directeurs et les quatre dimensions.
    \item Situer ITIL version~5 (2026) par rapport à ITIL~4.
\end{itemize}
\end{boxobjectifs}

\bigskip

% ============================================================================
\section*{Question directrice : que veut vraiment le client~?}
% ============================================================================

Chez KRIBILOG SA, le centre de services traite 4~800~tickets par mois, et
30~\% d'entre eux sont rouverts. Les indicateurs de l'infogérant NETPORT
SERVICES sont pourtant au vert~: les serveurs sont disponibles à
99,8~\%. Le directeur des opérations portuaires résume le malentendu~:
« Vos serveurs fonctionnent, mais mes conteneurs restent bloqués. »

Ce malentendu est l'objet d'ITIL~4. La Partie~\ref{part:pratique}, qui
s'ouvre ici, passe de la gouvernance à la pratique quotidienne des
services. Ce chapitre en pose les fondations~: l'histoire d'ITIL
(section~\ref{sec:c12-histoire}), ses concepts clés
(section~\ref{sec:c12-concepts}), le système de valeur des services et sa
chaîne de valeur (sections~\ref{sec:c12-svs} et~\ref{sec:c12-chaine}), les
quatre dimensions et les sept principes (sections~\ref{sec:c12-dimensions}
et~\ref{sec:c12-principes}), et enfin la place d'ITIL version~5
(section~\ref{sec:c12-itil5}).

\separateur

% ============================================================================
\section{Trente-Cinq Ans d'ITIL}
\label{sec:c12-histoire}
% ============================================================================

\begin{tableaucours}{L{2.4cm}L{4.0cm}L{5.8cm}}{%
    \tableauHeader{Version} &
    \tableauHeader{Logique} &
    \tableauHeader{Apport}
}
Fin des années 1980 & Recueil de pratiques & L'administration britannique
rassemble les bonnes pratiques de gestion de ses services informatiques \\
ITIL v2 (2000--2001) & Processus & Deux ouvrages centraux~: le soutien et la
fourniture des services~; diffusion mondiale \\
ITIL v3 (2007, révisé en 2011) & Cycle de vie & Cinq étapes~: stratégie,
conception, transition, exploitation et amélioration continue des
services \\
ITIL 4 (2019) & Système de valeur & Co-création de la valeur, 34~pratiques,
intégration de l'agilité, de DevOps et du \textit{lean} \\
ITIL version 5 (2026) & Produits et services & Voir la
section~\ref{sec:c12-itil5} \\
\end{tableaucours}
\vspace{-0.8em}
\begin{center}{\small\itshape Tableau~12.1. Évolution d'ITIL}\end{center}

Le passage d'ITIL~v3 à ITIL~4 est le plus profond. ITIL~v3 décrivait
\emph{comment le fournisseur gère ses services} au long de leur cycle de
vie. ITIL~4 décrit \emph{comment la valeur se crée entre le fournisseur et
ses clients}, et intègre les méthodes agiles et DevOps qui s'étaient
développées en dehors d'ITIL. Le mot « processus » cède la place au mot
« pratique », plus large~: une pratique comprend des processus, mais aussi
des personnes, des informations et des outils, comme les composantes de
COBIT (Chapitre~\ref{chap:cobit-archi}).

% ============================================================================
\section{Les Concepts Clés}
\label{sec:c12-concepts}
% ============================================================================

\subsection{Service et valeur}

\begin{boxdef}
\termecle{Service}\index{service (ITIL)}~: moyen de permettre la
co-création de valeur, en facilitant l'obtention des résultats que le
client veut atteindre, sans qu'il ait à gérer lui-même certains coûts et
certains risques.

\medskip
\termecle{Valeur}\index{valeur (ITIL)}~: bénéfices, utilité et importance
perçus de quelque chose. La valeur est \emph{perçue} par celui qui la
reçoit~: elle ne se décrète pas chez le fournisseur.

\medskip
\termecle{Co-création de valeur}\index{co-création de valeur}~: idée selon
laquelle la valeur naît de la collaboration entre le fournisseur et le
consommateur du service, et non de la seule action du fournisseur.
\end{boxdef}
{\small\itshape Définitions reformulées d'après ITIL~4~\cite{axelos2019itil4}.}

\subsection{Résultats, produits et relations}

ITIL~4 distingue soigneusement ce que le fournisseur \emph{produit} de ce
que le client \emph{obtient}.

\begin{tableaucours}{L{3.2cm}L{5.0cm}L{4.0cm}}{%
    \tableauHeader{Notion} &
    \tableauHeader{Sens} &
    \tableauHeader{Exemple KRIBILOG}
}
Livrable (\textit{output}) & Ce qu'une activité produit & Un écran de suivi
affiché, un ticket clos \\
Résultat (\textit{outcome}) & Ce que le client obtient grâce aux livrables
& Un conteneur localisé et chargé à l'heure \\
Produit & Configuration de ressources conçue pour offrir de la valeur &
Le logiciel de suivi, l'infrastructure, l'équipe de support \\
Offre de service & Description d'un ou plusieurs services, pour un groupe
de clients & « Suivi des conteneurs 24~h/24, support en 30~minutes » \\
Relation de service & Coopération entre fournisseur et consommateur &
Fourniture, consommation et gestion du service \\
\end{tableaucours}
\vspace{-0.8em}
\begin{center}{\small\itshape Tableau~12.2. Les notions clés d'ITIL~4
(reformulées)}\end{center}

Le malentendu de KRIBILOG s'éclaire~: NETPORT SERVICES mesure des
\emph{livrables} (serveurs disponibles, tickets clos), le directeur des
opérations attend un \emph{résultat} (conteneurs chargés). Un ticket clos
qui est rouvert le lendemain est un livrable sans résultat.

\subsection{Utilité et garantie}

\begin{boxdef}
\termecle{Utilité}\index{utilité} (\textit{utility})~: fonctionnalité
offerte par un service pour répondre à un besoin~: ce que fait le
service.

\medskip
\termecle{Garantie}\index{garantie} (\textit{warranty})~: assurance que le
service répondra aux exigences convenues~: disponibilité, capacité,
sécurité, continuité. C'est \emph{comment} le service le fait.
\end{boxdef}

Un service n'a de valeur que s'il a les deux~: une utilité sans garantie
(un logiciel parfait, mais indisponible un jour sur deux) n'en a pas
davantage qu'une garantie sans utilité.

% ============================================================================
\section{Le Système de Valeur des Services}
\label{sec:c12-svs}
% ============================================================================

\subsection{Cinq composantes}

Le cœur d'ITIL~4 est le \termecle{système de valeur des services}\index{système
de valeur des services} (\textit{service value system}, \sigle{SVS}). Il
décrit comment les composantes et les activités d'une organisation
fonctionnent ensemble pour transformer une \emph{opportunité} ou une
\emph{demande} en \emph{valeur}.

\noindent\begin{minipage}{\linewidth}
\begin{center}
\begin{tikzpicture}[
    entree/.style = {
        draw=CouleurAccent, fill=CouleurAccent!15, rounded corners=4pt,
        text width=2.2cm, minimum height=1.0cm, align=center,
        font=\footnotesize\bfseries\color{CouleurPrimaire}
    },
    comp/.style = {
        draw=CouleurBordDef, fill=CouleurDef, rounded corners=3pt,
        text width=3.6cm, minimum height=0.6cm, align=center,
        font=\scriptsize\bfseries\color{CouleurPrimaire}
    },
    cadre/.style = {
        draw=CouleurPrimaire, thick, rounded corners=6pt
    },
    fleche/.style = {-{Stealth[length=5pt]}, thick, color=CouleurSecondaire}
]
\node[entree] (op) at (0,0)   {Opportunité\\ et demande};
\node[comp] (p) at (4.2,1.3)  {Principes directeurs};
\node[comp] (g) at (4.2,0.6)  {Gouvernance};
\node[comp, fill=CouleurPrimaire, text=white] (c) at (4.2,-0.1)
    {Chaîne de valeur des services};
\node[comp] (pr) at (4.2,-0.8) {Pratiques};
\node[comp] (am) at (4.2,-1.5) {Amélioration continue};
\node[cadre, fit=(p)(am), inner sep=6pt] (svs) {};
\node[font=\scriptsize\bfseries\color{CouleurAccent}, above=0.05cm of svs]
    {Système de valeur des services};
\node[entree] (va) at (8.4,0) {Valeur};
\draw[fleche] (op) -- (svs.west |- op);
\draw[fleche] (svs.east |- va) -- (va);
\end{tikzpicture}
\end{center}
\vspace{-0.4em}
\begin{center}{\small\itshape Figure~12.1. Le système de valeur des
services d'ITIL~4 (représentation simplifiée)}\end{center}
\end{minipage}

\medskip

\begin{tableaucours}{L{3.6cm}L{8.6cm}}{%
    \tableauHeader{Composante} &
    \tableauHeader{Rôle}
}
Principes directeurs & Recommandations qui guident les décisions en toutes
circonstances (section~\ref{sec:c12-principes}) \\
Gouvernance & Moyens par lesquels l'organisation est dirigée et
contrôlée~: évaluer, diriger, surveiller \\
Chaîne de valeur des services & Modèle des six activités qui répondent à
la demande et créent la valeur (section~\ref{sec:c12-chaine}) \\
Pratiques & Ensembles de ressources organisationnelles conçus pour
accomplir un travail ou atteindre un objectif (34~pratiques) \\
Amélioration continue & Activité récurrente, à tous les niveaux, qui
assure que la performance répond aux attentes \\
\end{tableaucours}
\vspace{-0.8em}
\begin{center}{\small\itshape Tableau~12.3. Les composantes du système de
valeur des services}\end{center}

\subsection{ITIL et la gouvernance}

La deuxième composante est la \emph{gouvernance}. ITIL~4 y reprend le
modèle de l'ISO/IEC~38500~: évaluer, diriger, surveiller
(Chapitre~\ref{chap:normes}). C'est le point d'ancrage d'ITIL dans la
gouvernance des SI~: ITIL ne prétend pas organiser la gouvernance, mais il
s'y soumet. Le système de valeur des services d'une organisation doit être
dirigé et contrôlé par son organe de gouvernance, celui que COBIT organise
avec le domaine EDM.

% ============================================================================
\section{La Chaîne de Valeur des Services}
\label{sec:c12-chaine}
% ============================================================================

\subsection{Six activités}

\begin{tableaucours}{L{3.0cm}L{9.2cm}}{%
    \tableauHeader{Activité} &
    \tableauHeader{Finalité}
}
Planifier & Partager une vision commune de l'état actuel, des orientations
et des améliorations, pour tous les produits et services \\
Améliorer & Améliorer en continu les produits, les services et les
pratiques \\
S'engager & Comprendre les besoins des parties prenantes, entretenir la
transparence et de bonnes relations \\
Concevoir et faire la transition & Faire en sorte que les produits et
services répondent aux attentes de qualité, de coût et de délai \\
Obtenir ou construire & Disposer des composants de service nécessaires,
quand et où il le faut \\
Fournir et soutenir & Livrer et soutenir les services selon les
spécifications et les attentes convenues \\
\end{tableaucours}
\vspace{-0.8em}
\begin{center}{\small\itshape Tableau~12.4. Les six activités de la chaîne
de valeur des services (finalités reformulées)}\end{center}

\subsection{Des chaînes de valeur, pas un ordre fixe}

Les six activités ne forment pas un cycle à parcourir dans l'ordre, comme
les étapes d'ITIL~v3. Elles se combinent, dans un ordre propre à chaque
situation, en \termecle{flux de valeur}\index{flux de valeur}
(\textit{value streams})~: la suite d'étapes qu'une organisation enchaîne
pour créer un résultat donné pour un client. Une même pratique peut
intervenir dans plusieurs activités.

\begin{boxscenario}{Opération KRIBI. Conteneur bloqué}
\textbf{Déclencheur.} Un agent du port signale qu'un conteneur n'apparaît
plus dans le logiciel de suivi.

\medskip
\textbf{Flux de valeur.}
\begin{enumerate}
    \item \emph{S'engager}~: le centre de services enregistre l'appel et
          comprend l'impact (un navire attend).
    \item \emph{Fournir et soutenir}~: le support diagnostique, contourne
          le problème et relocalise le conteneur.
    \item \emph{Obtenir ou construire}~: l'éditeur corrige l'anomalie
          logicielle.
    \item \emph{Concevoir et faire la transition}~: la correction est
          testée, autorisée, puis déployée.
    \item \emph{Améliorer}~: l'analyse de cause alimente le registre des
          problèmes.
\end{enumerate}

\textbf{Constatation.} Chez KRIBILOG, le flux s'arrête aujourd'hui à
l'étape~2~: le ticket est clos dès le contournement, sans correction de
la cause. Il est rouvert lors du blocage suivant.

\medskip
\textbf{Analyse} (confiance élevée). Le taux de réouverture de 30~\%
s'explique par l'absence des étapes~3 à~5 dans le flux réel~; les
pratiques de gestion des problèmes et d'habilitation des changements
(Chapitres~\ref{chap:itil-incidents} et~\ref{chap:itil-changements}) sont
les premières à mettre en place.
\end{boxscenario}

% ============================================================================
\section{Les Quatre Dimensions}
\label{sec:c12-dimensions}
% ============================================================================

ITIL~4 rappelle qu'un service ne se réduit pas à sa technique. Toute
démarche doit tenir compte de quatre dimensions, sous l'influence de
facteurs externes (politiques, économiques, sociaux, technologiques,
juridiques et environnementaux).

\begin{tableaucours}{L{3.6cm}L{4.4cm}L{4.2cm}}{%
    \tableauHeader{Dimension} &
    \tableauHeader{Contenu} &
    \tableauHeader{Chez KRIBILOG}
}
Organisations et personnes & Structures, rôles, compétences, culture &
Centre de services sous-dimensionné, pas de responsable des problèmes \\
Information et technologie & Données, applications, infrastructures,
outils & Logiciel de suivi, outil de tickets, base de connaissances
absente \\
Partenaires et fournisseurs & Relations avec les autres organisations qui
contribuent au service & NETPORT SERVICES, l'éditeur du logiciel \\
Flux de valeur et processus & Enchaînement des activités et des
processus & Flux « conteneur bloqué » interrompu \\
\end{tableaucours}
\vspace{-0.8em}
\begin{center}{\small\itshape Tableau~12.5. Les quatre dimensions de la
gestion des services}\end{center}

\begin{boxalert}
\textbf{Violation typique.} Un projet ITIL est réduit à l'achat d'un outil
de gestion des tickets. L'outil est installé, mais les rôles ne sont pas
définis, l'infogérant n'est pas tenu par le contrat d'y saisir ses
interventions, et aucun flux de valeur n'est décrit.

\smallskip
\textbf{Règle~:} une démarche ITIL traite toujours les quatre dimensions.
Un outil ne couvre qu'une partie de la deuxième.
\end{boxalert}

% ============================================================================
\section{Les Sept Principes Directeurs}
\label{sec:c12-principes}
% ============================================================================

\begin{tableaucours}{L{3.8cm}L{8.4cm}}{%
    \tableauHeader{Principe} &
    \tableauHeader{Application}
}
Privilégier la valeur & Tout ce que l'on fait doit, directement ou
indirectement, créer de la valeur pour les parties prenantes \\
Commencer là où l'on est & Partir de l'existant, l'observer et le
mesurer, plutôt que de tout reconstruire \\
Progresser par itérations, avec retour d'information & Avancer par
étapes courtes, en recueillant les retours à chaque étape \\
Collaborer et promouvoir la visibilité & Associer les bonnes personnes,
rendre le travail et ses résultats visibles \\
Penser et travailler de manière globale & Traiter le service comme un
tout, avec ses quatre dimensions \\
Rester simple et pratique & Supprimer ce qui ne crée pas de valeur,
préférer le nombre minimal d'étapes \\
Optimiser et automatiser & Optimiser d'abord, automatiser ensuite~; ne pas
automatiser un processus défaillant \\
\end{tableaucours}
\vspace{-0.8em}
\begin{center}{\small\itshape Tableau~12.6. Les sept principes directeurs
d'ITIL~4 (notre traduction)}\end{center}

\begin{boxmetier}{Responsable des services IT}
\textbf{Appliquer les principes au projet ITIL de KRIBILOG}
\begin{enumerate}
    \item \emph{Privilégier la valeur}~: mesurer les conteneurs bloqués et
          les délais de chargement, pas seulement la disponibilité des
          serveurs.
    \item \emph{Commencer là où l'on est}~: analyser les 4~800~tickets du
          dernier mois avant de choisir un nouvel outil.
    \item \emph{Progresser par itérations}~: traiter d'abord la gestion des
          problèmes sur le logiciel de suivi, puis étendre.
    \item \emph{Collaborer}~: associer les agents du port et NETPORT
          SERVICES à la définition des priorités.
    \item \emph{Optimiser et automatiser}~: n'automatiser l'affectation des
          tickets qu'une fois la classification stabilisée.
\end{enumerate}
\end{boxmetier}

\subsection{Les 34 pratiques}

ITIL~4 décrit 34~pratiques, réparties en trois familles~\cite{axelos2019itil4}.

\begin{tableaucours}{L{3.2cm}cL{8.0cm}}{%
    \tableauHeader{Famille} &
    \tableauHeader{Nb} &
    \tableauHeader{Exemples}
}
Management général & 14 & Amélioration continue, gestion des risques,
gestion de la sécurité de l'information, gestion des fournisseurs, gestion
des projets, gestion du portefeuille, gestion des relations \\
Gestion des services & 17 & Centre de services, gestion des incidents,
des problèmes, des demandes de service, habilitation des changements,
gestion des mises en production, des niveaux de service, de la
disponibilité, de la continuité, de la configuration \\
Gestion technique & 3 & Gestion des déploiements, gestion des
infrastructures et des plateformes, développement et gestion des
logiciels \\
\end{tableaucours}
\vspace{-0.8em}
\begin{center}{\small\itshape Tableau~12.7. Les trois familles de pratiques
d'ITIL~4}\end{center}

Les Chapitres~\ref{chap:itil-incidents} à~\ref{chap:itil-niveaux}
présentent les pratiques de gestion des services les plus importantes pour
la gouvernance~; le Chapitre~\ref{chap:processus} les met en
correspondance avec COBIT et l'ISO/IEC~20000-1.

\begin{boxobjet}{Amélioration continue (pratique ITIL 4)}
\textbf{Finalité~:} aligner en permanence les pratiques et les services
sur les besoins de l'organisation, par l'amélioration des produits, des
services et des pratiques.

\textbf{Modèle en sept étapes~:} quelle est la vision~? Où en
sommes-nous~? Où voulons-nous aller~? Comment y arriver~? Agir. Y
sommes-nous arrivés~? Comment garder l'élan~?

\textbf{Sorties~:} registre des améliorations, plans d'action, mesures
avant et après.

\textbf{RACI (extrait, KRIBILOG)~:} \emph{A} au responsable des services
informatiques, \emph{R} aux propriétaires des pratiques, \emph{C} aux
métiers et à l'infogérant.

\textbf{Indicateurs~:} nombre d'améliorations menées à terme, gain mesuré
par amélioration.

\textbf{Correspondance COBIT~:} objectifs APO11 (qualité) et MEA01
(surveillance de la performance).
\end{boxobjet}

On remarquera que le modèle d'amélioration continue d'ITIL reprend les
sept questions de la feuille de route de COBIT~2019
(Chapitre~\ref{chap:cobit-capacite})~: les deux référentiels partagent la
même logique d'amélioration.

% ============================================================================
\section{Et ITIL Version 5~?}
\label{sec:c12-itil5}
% ============================================================================

ITIL version~5 a été annoncée par PeopleCert en janvier~2026, et ses
modules de formation sont publiés progressivement depuis~\cite{peoplecert2026itil5}.
Les sources publiques disponibles permettent d'établir les points
suivants.

\begin{tableaucours}{L{3.4cm}L{4.2cm}L{4.6cm}}{%
    \tableauHeader{Aspect} &
    \tableauHeader{ITIL 4 (2019)} &
    \tableauHeader{ITIL version 5 (2026)}
}
Objet & Gestion des services & Produits et services numériques, gérés
ensemble \\
Modèle central & Système et chaîne de valeur des services & Cycle de vie
des produits et des services, de bout en bout \\
Principes directeurs & Sept principes & Les mêmes sept principes, selon
les présentations disponibles \\
Pratiques & 34~pratiques en trois familles & 34~pratiques, regroupées
autrement selon les présentations disponibles \\
Intelligence artificielle & Peu traitée & Présentée comme intégrée dès la
conception, avec un module consacré à la gouvernance de l'IA \\
Certifications & Toujours reconnues & Nouveau schéma, introduit
progressivement \\
\end{tableaucours}
\vspace{-0.8em}
\begin{center}{\small\itshape Tableau~12.8. ITIL~4 et ITIL version~5
(d'après les sources publiques disponibles en octobre~2026)}\end{center}

\begin{boiteRemarque}{Pourquoi ce cours reste-t-il sur ITIL~4~?}
Trois raisons. D'abord, ITIL~4 est la version que pratiquent aujourd'hui
les organisations où exerceront les étudiants. Ensuite, ses concepts sont
largement repris par la version~5~: principes, pratiques, co-création de
la valeur. Enfin, le modèle central de la version~5 n'est pas encore
documenté en détail dans les sources publiques~: l'enseigner exigerait de
s'appuyer sur la publication officielle. Les étudiants qui passeront une
certification devront vérifier la version exigée par l'examen.
\end{boiteRemarque}

% ============================================================================
\separateur
% ============================================================================

\begin{boxresume}
\textbf{Ce qu'il faut retenir de ce chapitre~:}
\begin{itemize}
    \item \textbf{De processus en système de valeur}~: ITIL est passé d'un
          recueil de processus (v2) à un cycle de vie (v3), puis à un
          système de valeur des services (ITIL~4, 2019).
    \item \textbf{La valeur est co-créée}~: elle est perçue par le client.
          Ne pas confondre livrable (ce que l'on produit) et résultat (ce
          que le client obtient).
    \item \textbf{Utilité et garantie}~: ce que fait le service, et
          l'assurance qu'il le fera comme convenu.
    \item \textbf{Cinq composantes}~: principes directeurs, gouvernance,
          chaîne de valeur, pratiques, amélioration continue. La
          gouvernance reprend le modèle évaluer, diriger, surveiller.
    \item \textbf{Six activités} qui se combinent en flux de valeur, sans
          ordre imposé.
    \item \textbf{Quatre dimensions et sept principes}~: un projet ITIL n'est
          jamais un simple projet d'outil.
    \item \textbf{ITIL version~5} (2026) recentre le référentiel sur les
          produits et services numériques~; ce cours reste sur ITIL~4.
\end{itemize}
\end{boxresume}

\bigskip

\begin{boxexo}[title={\ding{45}\quad Exercice~12.1. Livrable ou résultat~? \niveaubase}]
Pour chacun des éléments suivants, indiquez s'il s'agit d'un livrable ou
d'un résultat, et précisez pour quel client.
\begin{enumerate}[(a)]
    \item Un ticket clos dans le délai.
    \item Un fournisseur de l'État payé à bonne date par l'ANPP.
    \item Un serveur disponible à 99,8~\%.
    \item Un client de la BANQUE SANAGA qui retire de l'argent un dimanche.
    \item Un rapport mensuel de NETPORT SERVICES.
\end{enumerate}
\end{boxexo}

\begin{boxexo}[title={\ding{45}\quad Exercice~12.2. Un flux de valeur à la BANQUE SANAGA \niveaumoyen}]
Décrivez le flux de valeur « ouverture d'un compte par l'application
mobile », de la demande du client jusqu'au premier virement.
\begin{enumerate}[(a)]
    \item Identifiez les étapes et rattachez chacune à une activité de la
          chaîne de valeur.
    \item Pour chaque étape, indiquez la pratique ITIL concernée.
    \item Identifiez, pour chacune des quatre dimensions, un risque qui
          pourrait interrompre le flux.
\end{enumerate}
\end{boxexo}

\begin{boxexo}[title={\ding{45}\quad Exercice~12.3. Le contrat de KRIBILOG au prisme d'ITIL~4 \niveauexpert}]
Le contrat entre KRIBILOG et NETPORT SERVICES fixe un seul engagement~: la
disponibilité des serveurs à 99,5~\%.
\begin{enumerate}[(a)]
    \item Expliquez, avec les notions d'utilité, de garantie, de livrable et
          de résultat, pourquoi cet engagement ne protège pas KRIBILOG.
    \item Proposez trois engagements orientés vers le résultat, mesurables.
    \item Montrez comment ces engagements s'inscrivent dans la gouvernance
          de KRIBILOG (qui les fixe, qui les surveille, qui en rend compte).
    \item Ce travail sera repris au Chapitre~\ref{chap:itil-niveaux}~:
          indiquez les pratiques ITIL qu'il mobilisera.
\end{enumerate}
\end{boxexo}

\bigskip

\begin{boxinfo}
\textbf{Transition.} Le système de valeur fixe le cadre. Le
Chapitre~\ref{chap:itil-incidents} entre dans les pratiques du quotidien~:
le centre de services, la gestion des incidents, des problèmes et des
demandes de service, là où se joue la perception du service par ses
utilisateurs.
\end{boxinfo}
